% Generated draft: requires mathematical and source review.
\documentclass[11pt]{article}
\usepackage[margin=25mm]{geometry}
\usepackage{fontspec}
\setmainfont{Latin Modern Roman}
\usepackage{amsmath,amssymb,mathtools,unicode-math}
\setmathfont{Latin Modern Math}
\usepackage{enumitem,needspace,xurl,hyperref}
\setlength{\parindent}{0pt}
\setlength{\parskip}{5pt}
\setlength{\emergencystretch}{3em}
\newcommand{\N}{\mathbb{N}}
\newcommand{\Z}{\mathbb{Z}}
\newcommand{\Q}{\mathbb{Q}}
\newcommand{\R}{\mathbb{R}}
\newcommand{\C}{\mathbb{C}}
\newcommand{\F}{\mathbb{F}}
\newcommand{\A}{\mathbb{A}}
\newcommand{\PP}{\mathbb{P}}
\newcommand{\E}{\mathbb{E}}
\newcommand{\eps}{\varepsilon}
\newcommand{\dd}{\,\mathrm{d}}
\begin{document}
\section*{UnsolvedMath ID: 800724}
\textbf{Source problem code:} catalogue entry
\subsection{Definitions and mathematical statement}
Let \(F\) be a field and let \(A\) be a finite-dimensional central simple \(F\)-algebra. These ambient hypotheses are inferred from the source title and notation; the supplied excerpt does not state them explicitly. For each field extension \(L/F\), put \(A_L=A\otimes_F L\). For a central simple algebra \(B\) over a field \(K\), let \(B^\times\) be its multiplicative group, let \(\operatorname{Nrd}_{B/K}:B^\times\to K^\times\) be its reduced norm, and define
\[
\operatorname{SL}_1(B)=\ker(\operatorname{Nrd}_{B/K}),\qquad
\operatorname{SK}_1(B)=\operatorname{SL}_1(B)/[B^\times,B^\times],
\]
where \([B^\times,B^\times]\) is the subgroup generated by multiplicative commutators. The equation \(\operatorname{SK}_1(B)=0\) means that this quotient group is trivial.

Write \(A\cong M_m(D)\), where \(D\) is a central division \(F\)-algebra, and define \(\operatorname{ind}(A)=\sqrt{\dim_F D}\). A positive integer is square-free if no square of a prime divides it.

\textbf{Conjecture 39 (attributed in the source to Suslin).} For every field \(F\) and every finite-dimensional central simple \(F\)-algebra \(A\),
\[
\left(\operatorname{SK}_1(A_L)=0\ \text{for every field extension }L/F\right)
\quad\Longrightarrow\quad
\operatorname{ind}(A)\text{ is square-free}.
\]
The quantifier over extensions has no stated restriction, so it includes arbitrary field extensions.

The source describes this conjecture as a converse to Wang's result. Its preceding paragraph discusses an unspecified theorem: it attributes the ``if'' direction to [Vos, p.~186], the full theorem to that reference together with [Gille:KT, Theorem~5.9], and an additional equivalence to the abstract group \(\operatorname{SL}_r(A)(L)\) being ``projectively simple'' for every extension \(L/F\) and every integer \(r\geq 2\). The theorem's conditions and the meaning of ``projectively simple'' are absent from the supplied excerpt; consequently, that contextual equivalence cannot be stated precisely here.
\subsection{Short English statement}
Suslin's conjecture asks whether a central simple algebra must have square-free index whenever its reduced Whitehead group becomes trivial over every field extension of its center. The ambient algebra hypotheses are inferred from context. The registry labels the problem open; this statement-only draft does not independently verify that status.
\subsection{Sources}
\begin{itemize}
\item Local UnsolvedMath registry, permanent ID \(800724\), source code \texttt{catalogue entry}, title: \emph{Open Problems on Central Simple Algebras --- Conjecture 39}. Supplied upstream revision: \texttt{d371ee917ac9647fbb06cd6fc529deec6f3dda82}. The record lists Auel, Brussel, Garibaldi, and Vishne as proposers, gives the year 2010, and labels the problem open.
\item Auel, Brussel, Garibaldi, and Vishne, \emph{Open Problems on Central Simple Algebras} (2010), Conjecture 39: \url{https://arxiv.org/abs/1006.3304}. The supplied provenance records access on 2026-10-16 and extraction by \texttt{arxiv\_latex}; these are supplied provenance claims, not independent source verification.
\item The conjecture is attributed in the supplied excerpt to Suslin, reference [Suslin:SK1, p.~75]. Full bibliographic details for this reference are not supplied.
\item The contextual paragraph cites [Vos, p.~186] and [Gille:KT, Theorem~5.9]. Full bibliographic details and the preceding theorem are not supplied. The excerpt also mentions Wang's result without a bibliographic citation.
\end{itemize}
\end{document}
